\section{Introduction}
\label{sec:introduction}

Modern organisations generate enormous volumes of security-relevant data:
authentication events, process executions, network connections, file changes,
and system errors. Making sense of this data --- and acting on it quickly
enough to matter --- is the work of a Security Operations Centre (SOC). At the
centre of every SOC sits a Security Information and Event Management (SIEM)
platform, which collects logs from across an estate, correlates them, and
raises alerts an analyst can investigate.

For someone entering the cybersecurity profession, this creates a practical
difficulty. Almost every junior SOC and security-analyst role expects
familiarity with a SIEM, yet undergraduate study tends to teach the theory of
detection and incident response without sustained hands-on exposure to the
tools themselves. A graduate can describe what a SIEM does without ever having
operated one.

This project was undertaken to close that gap directly. It designs, builds,
and documents a working SOC home lab using only free and open-source software,
producing both a functional environment and a reproducible record of how it
was created.

\subsection{Aim and Objectives}
\label{sec:aim}

The \textbf{aim} of this project is to design, deploy, and operate a functional
SOC home lab built on an open-source SIEM, capable of collecting and analysing
security telemetry from multiple endpoints and supporting a realistic analyst
triage workflow.

The aim is met through the following \textbf{objectives}:

\begin{enumerate}[label=\textbf{O\arabic*.},leftmargin=*]
    \item Deploy an open-source SIEM server on an isolated virtual network.
    \item Configure two monitored endpoints --- a Windows host and a Linux
          host --- each reporting to the SIEM.
    \item Enable enhanced host telemetry beyond default operating-system logs.
    \item Verify centralised log collection and establish a baseline of normal
          activity.
    \item Simulate a set of security events and observe the resulting alerts.
    \item Triage the alerts, classifying them and documenting the reasoning.
    \item Produce a fully reproducible record of the build so the lab can be
          rebuilt by a third party.
\end{enumerate}

\subsection{Research Questions}
\label{sec:rq}

The project is framed around three research questions, revisited in the
Discussion and Conclusions:

\begin{itemize}[leftmargin=*]
    \item \textbf{RQ1.} Can a credible SOC monitoring environment, capable of
          collecting endpoint telemetry and generating meaningful alerts, be
          built using only free and open-source tools?
    \item \textbf{RQ2.} What additional configuration is required beyond a
          default SIEM and operating-system installation to achieve
          detection-grade visibility?
    \item \textbf{RQ3.} To what extent does a single-machine virtualised lab
          realistically represent the work of a professional SOC, and where
          are the limits of that representation?
\end{itemize}

\subsection{Contributions}
\label{sec:contributions}

This project contributes:

\begin{itemize}[leftmargin=*]
    \item A complete, working SOC home-lab build using Wazuh, deployed on an
          isolated network with Windows and Linux endpoints.
    \item A fully documented, step-by-step record enabling independent
          reproduction of the environment.
    \item A set of attack simulations paired with the alerts they produce and
          an analyst-style triage of each.
    \item An honest assessment of how closely such a lab represents real SOC
          work, including its limitations.
\end{itemize}

\subsection{Report Structure}
\label{sec:structure}

The remainder of this report is organised as follows.
\textbf{Section~\ref{sec:litreview}} reviews the background literature on
SOCs, SIEM technology, logging, and security home labs, and positions this
work against related projects.
\textbf{Section~\ref{sec:methodology}} sets out the general methodology and
design decisions.
\textbf{Section~\ref{sec:implementation}} documents the specific
implementation of the lab.
\textbf{Section~\ref{sec:evaluation}} evaluates whether the lab meets its
objectives.
\textbf{Section~\ref{sec:discussion}} discusses challenges and revisits the
research questions, and \textbf{Section~\ref{sec:conclusions}} concludes and
identifies future work.
